\documentclass{article}
\usepackage{amsmath,hyperref}
\title{Choose the successor at the first charged effect}
\author{}
\date{}
\begin{document}
\maketitle
\paragraph{Admitted family and invariant.}
One old duty is owed and either of two named candidates can settle it
once. An independent obligation must be preserved. After an initial
proposal, a native plant event can change the set of usable candidates.
At the service cut there is time or capacity for one charged action
in the tight cases. A bounded local adapter has authority to select
either candidate and to query feasibility without consuming the action.
No competing mutator can change the checked premises from that read
through the committed effect. This is one serialized source loop in
the fixtures, not a distributed lock or an assumption that any remote
controller can inspect global state. Neither adapter has to predict
which change occurs before the service cut.

Run the same finite selection routine at that cut. Under a source
contract saying each positive read really supports settlement, the
routine dispatches only a supported action, at most once. If there is a
supported candidate it makes one commit and preserves the unrelated
work already settled by then. Otherwise it refuses this duty without
spending the action. The Lean proof establishes finding a supported
candidate whenever one exists, and that no selected candidate is
unsupported. The argument from the single call site to one actual
effect depends on the adapter's serialization and correspondence
premises. This is the familiar read-and-commit discipline
associated with a shared object
(\href{https://www.cs.cmu.edu/~wing/publications/HerlihyWing90.pdf}{Herlihy and Wing}),
not a novel concurrency theorem.

\paragraph{Two native adapters.}
In the pinned RWARE model, one of two cells can be occupied by an old
carrier. The initial proposal points to the other. Before the service
cut a native move can change which is occupied; one can also have
both occupied. Each of the two successors can enter its own cell in
one native invocation. One unrelated robot completes an actual
move without yielding that progress to the repair. The selector uses
the live occupancy grid and the existing native move; it does not
time-step each candidate just to test admission. A blind blocked
move still consumes the service tick, and is refused by the original
collision rules. The service authorization and equivalence of the
two target cells for the old duty are fixture assumptions, not a
stock dispatch contract.

In the pinned SubT broker, an independent source really sends a
payload and incurs a receiver debit. A subsequent independent
delivery can put live debit on the chosen receiver (or both). The
exact previous receive row remains active at a gap of exactly
one second; it has expired at $17/16$. The adapter reads live
local broker debits and known configured cap at the same serialized
cut as dispatch, then chooses where to send a four-byte old report.
A failed receive does not imply no effect: the broker already debited
the sender before testing receive capacity. At a cap of four, retrying
after the wrong choice cannot deliver the old report, even when the
other receiver was usable all along. At the exact one-second boundary,
both receivers can be over budget even when the crossed background
traffic would leave one usable $1/16$ later. The replay records both
background deliveries. It uses the inherited fixed high-power and
fixed random-draw fixture, not a field assurance or claim that a
stock remote SubT client can inspect these debits.

\paragraph{Comparison.}
A preflight is cheap but can go stale; the same preflight binds opposite
late choices. Neither fixed order guarantees service in the tight
cut across the two possible later states. In the radio case an eight-byte
cap can instead let a charged follow-up work, when only the first
recipient is blocked. A direct scheduler with the same co-located,
serialized inspection and the same single commit agrees on every
case in both adapters. Noninterference is a mechanism requirement:
a check made just before a competing mutator can be as stale as a
preflight, and the stock implementations were not shown to export a
cross-process atomic selector. Reusing the generic selector removes a
specialized implementation, not a measured service or compute
advantage over equal-information direct scheduling. Action dependencies
can encode the handoff if the source check remains valid through use
(see \href{https://whoenig.github.io/publications/2019_RA-L_Hoenig.pdf}{H\"onig et al.}
for effect-notified execution and replanning). This is a bounded,
source-backed transfer across occupancy and budget, not a full I1
operating envelope, a source-grounded common-cause recovery, or an
independent I2 mission.
\end{document}
